\section{Определение матрицы линейного оператора. Формулировка теоремы о преобразовании матрицы линейного оператора при переходе от одного базиса к другому}

\begin{definition}[Матрица линейного оператора]
Пусть $V$ --- линейное пространство над полем $F$, $E = \{e_1, \ldots, e_n\}$ --- базис в $V$, $\varphi \in L(V, V)$.

Образы базисных векторов: $\varphi(e_1), \ldots, \varphi(e_n) \in V$.

Разложим их по базису $E$:
\begin{align*}
\varphi(e_1) &= a_{11}e_1 + \cdots + a_{n1}e_n, \\
&\vdots \\
\varphi(e_n) &= a_{1n}e_1 + \cdots + a_{nn}e_n.
\end{align*}

Матрица
\[
A_\varphi^E = \begin{pmatrix}
a_{11} & \cdots & a_{1n} \\
\vdots & \ddots & \vdots \\
a_{n1} & \cdots & a_{nn}
\end{pmatrix}
\]
называется \textbf{матрицей линейного оператора} $\varphi$ в базисе $E$. Столбцы матрицы --- это координаты векторов $\varphi(e_1), \ldots, \varphi(e_n)$.
\end{definition}

\begin{theorem}[Преобразование матрицы оператора при смене базиса]
Пусть $V, F, E, E'$ --- два базиса в $V$, $\varphi \in L(V, V)$. Тогда
\[
A_\varphi^{E'} = T^{-1} A_\varphi^E T,
\]
где $T$ --- матрица перехода от базиса $E$ к базису $E'$.
\end{theorem}
